% Bagian: first-order-logic
% Bab: introduction
% Subbagian: sentences

\documentclass[../../../include/open-logic-section]{subfiles}

\begin{document}

\olfileid[id]{fol}{int}{snt}

\olsection{\usetoken{P}{sentence}}

Baik, sekarang kita mempunyai (sketsa) definisi keterpenuhan (``benar
dalam'') bagi !!{structure}s dan !!{formula}s. Namun, definisi itu memerlukan
unsur tambahan ini---penugasan untuk !!a{variable}---sedangkan yang kita
inginkan ialah definisi !!{sentence}s. Bagaimana kita menyingkirkan penugasan,
dan apa itu !!{sentence}s?

Anda mungkin ingat pembahasan dalam pengantar pertama Anda tentang logika
formal mengenai hubungan antara !!{variable}s dan kuantor. Kuantor selalu
diikuti oleh !!a{variable}, lalu dalam bagian !!{sentence} yang dikenai
kuantor tersebut (``cakupan'' kuantor itu), kita memahami bahwa !!{variable}
tersebut ``terikat'' oleh kuantor itu. Dalam !!{formula}s, tidak diharuskan
bahwa setiap !!{variable} mempunyai kuantor yang bersesuaian, dan
!!{variable}s tanpa kuantor yang bersesuaian disebut ``bebas'' atau ``tak
terikat.'' Kita akan menganggap !!{sentence}s sebagai semua !!{formula}s yang
tidak mempunyai !!{variable}s bebas.

Sekali lagi, gagasan intuitif mengenai kapan suatu kemunculan !!a{variable}
dalam !!a{formula}~$!A$ terikat, kuantor mana yang mengikatnya, dan kapan
kemunculan itu bebas, tidaklah sulit dipahami. Anda mungkin telah mempelajari
suatu metode untuk mengujinya, barangkali dengan menghitung tanda kurung. Kita
harus menuntut suatu definisi yang tepat---dan karena kita telah
mendefinisikan !!{formula}s melalui induksi, kita juga dapat memberikan
definisi mengenai kemunculan bebas dan terikat suatu !!a{variable}~$x$ dalam
!!a{formula}~$!A$ melalui induksi. Misalnya, bagi bahasa kita yang
disederhanakan, definisi tersebut dapat berbentuk seperti berikut:
\begin{enumerate}
  \item Jika $!A$ atomik, semua kemunculan $x$ di dalamnya bebas (yakni,
  kemunculan $x$ dalam~$\Atom{\Obj P}{x}$ bebas).
  \item Jika $!A$ berbentuk $\lnot !B$, suatu kemunculan~$x$ dalam~$\lnot
  !B$ bebas jika dan hanya jika kemunculan~$x$ yang bersesuaian bebas
  dalam~$!B$ (yakni, kemunculan bebas variabel dalam~$!B$ tepat sama dengan
  kemunculan yang bersesuaian dalam~$\lnot !B$).
  \item Jika $!A$ berbentuk $(!B \land !C)$, suatu kemunculan~$x$
  dalam~$(!B \land !C)$ bebas jika dan hanya jika kemunculan~$x$ yang
  bersesuaian bebas dalam~$!B$ atau dalam~$!C$.
  \item Jika $!A$ berbentuk $\lexists[x][!B]$, tidak ada kemunculan $x$
  dalam~$!A$ yang bebas; jika formula tersebut berbentuk $\lexists[y][!B]$
  dengan $y$ suatu !!{variable} yang berbeda dari~$x$, suatu kemunculan~$x$ dalam
  $\lexists[y][!B]$ bebas jika dan hanya jika kemunculan~$x$ yang
  bersesuaian bebas dalam~$!B$.
\end{enumerate}

Setelah kita mempunyai definisi yang tepat mengenai kemunculan bebas dan
terikat variabel, kita cukup mengatakan: !!a{sentence} ialah setiap
!!{formula} tanpa kemunculan bebas !!{variable}s.

\end{document}
